\documentclass[11pt]{article}

\usepackage[T1]{fontenc}
\usepackage[utf8]{inputenc}
\usepackage{lmodern}
\usepackage{amsmath,amssymb,amsthm}
\usepackage[margin=1in]{geometry}
\usepackage{microtype}
\usepackage{hyperref}

\hypersetup{
  colorlinks=true,
  linkcolor=blue,
  citecolor=blue,
  urlcolor=blue,
  pdftitle={A limiting density for the distinct-variable Erdos Problem 302}
}

\newtheorem{theorem}{Theorem}[section]
\newtheorem{lemma}[theorem]{Lemma}
\newtheorem{proposition}[theorem]{Proposition}
\newtheorem{corollary}[theorem]{Corollary}
\theoremstyle{definition}
\newtheorem{definition}[theorem]{Definition}
\theoremstyle{remark}
\newtheorem{remark}[theorem]{Remark}

\numberwithin{equation}{section}

\newcommand{\N}{\mathbb N}
\newcommand{\eps}{\varepsilon}
\newcommand{\ind}{\mathbf 1}
\newcommand{\lcm}{\operatorname{lcm}}
\newcommand{\supp}{\operatorname{supp}}

\title{A limiting density for the distinct-variable\\ Erd\H{o}s Problem 302}
\author{Dmitry Khanukov}
\date{17 September 2026}

\begin{document}
\maketitle

\begin{abstract}
Let \(f(N)\) be the largest cardinality of a subset of
\(\{1,\ldots,N\}\) containing no pairwise distinct \(a,b,c\) with
\(1/a=1/b+1/c\). We prove that \(f(N)/N\) has a limit. The main
estimate is a uniform vertex cover for solutions whose larger primitive
parameter \(v\) exceeds a prescribed cutoff: its size is
\(O(N(\log R)^{-1/1000})\), for every \(N\). The proof combines a
simultaneous small-prime-factor exceptional-set estimate, a uniform
Selberg upper bound for arbitrary marked sets of primes, and an
upper-bound specialization of Kevin Ford's divisor-in-an-interval
theorem. Finite-prime decompositions then give an exact variational
characterization and quantitative finite truncation bounds for the limiting
density. The present argument does not determine the value of the density
and therefore does not completely solve Problem 302.
\end{abstract}

\section{Introduction and main statement}

We consider the distinct-variable extremal question posed by Erd\H{o}s and
Graham \cite{EG1980} and recorded as Erd\H{o}s Problem 302
\cite{EP302}. Throughout,
\(\N=\{1,2,\ldots\}\), \([N]=\{1,\ldots,N\}\) for \(N\in\N\), and
all logarithms are natural.

A set \(A\subseteq\N\) is \emph{reciprocal-triple-free} if it contains
no pairwise distinct positive integers satisfying
\begin{equation}\label{eq:reciprocal}
 \frac1a=\frac1b+\frac1c.
\end{equation}
In any such solution, \(a\) is the smallest entry; interchanging
\(b,c\) permits the ordering \(a<b<c\). Define
\[
 f(N)=\max\{|A|:A\subseteq[N]\text{ is reciprocal-triple-free}\}.
\]
The restriction to pairwise distinct entries is part of this
definition.

The problem page records the elementary lower bound \(5N/8+o(N)\)
and an upper bound \(9N/10+o(N)\) \cite{EP302}. A recent preliminary
preprint gives a qualitative improvement above \(5/8\) and the upper
constant \(140803024/163562355\) \cite{Khanukov2026}. Those bounds do
not imply convergence of \(f(N)/N\). The purpose of the present paper
is only to establish that convergence and to give the variational
description below.

Every ordered solution \(a<b<c\) has a unique representation
\begin{equation}\label{eq:param-intro}
 a=kuv,\qquad b=ku(u+v),\qquad c=kv(u+v),
 \qquad k\in\N,\quad 1\le u<v,\quad (u,v)=1;
\end{equation}
see Lemma~\ref{lem:param}. For an integer \(R\ge2\), let \(f_R(N)\)
be the largest size of a subset of \([N]\) avoiding only the
solutions \eqref{eq:param-intro} with \(v\le R\). Thus
\[
 f(N)\le f_{R+1}(N)\le f_R(N)\le N.
\]

\begin{theorem}\label{thm:main}
The following assertions hold.
\begin{enumerate}
\item For every integer \(R\ge2\), the limit
\[
 \alpha_R=\lim_{\substack{N\to\infty\\N\in\N}}\frac{f_R(N)}N
\]
exists.
\item There are absolute constants \(C>0\) and
\(R_0\in\N\), with \(R_0\ge2\), such that for every integer
\(R\ge R_0\) and every \(N\in\N\),
\begin{equation}\label{eq:uniform-comparison-main}
 0\le f_R(N)-f(N)\le CN(\log R)^{-1/1000}.
\end{equation}
\item The limit
\[
 \alpha=\lim_{\substack{N\to\infty\\N\in\N}}\frac{f(N)}N
\]
exists, and
\begin{equation}\label{eq:alpha-inf-main}
 \alpha=\inf_{\substack{R\ge2\\R\in\N}}\alpha_R,
 \qquad
 0\le\alpha_R-\alpha\le C(\log R)^{-1/1000}
 \quad(R\ge R_0).
\end{equation}
\item An exact variational characterization and quantitative finite
truncation bounds are given in Theorem~\ref{thm:effective}.
\end{enumerate}
\end{theorem}


The substantive approximation statement is
\eqref{eq:uniform-comparison-main}. We prove the stronger assertion
that all solutions with \(v>R\) have a vertex cover of the stated
size, independently of the choice of a truncated independent set.
For the analytic argument, large values of the scaling parameter
\(k\) are treated by a weighted sieve, while small values are treated
by Ford's theorem \cite[Theorem~1]{Ford2008}. The precise
divisor-interval range used from that theorem is stated in
Theorem~\ref{thm:ford}.

The Selberg construction used below is classical; see, for example,
\cite{HR}. We include its quadratic-form proof and all remainder
estimates because uniformity over arbitrary marked sets of primes
is essential. Ford's theorem supplies the divisor-distribution
input, not a result about the extremal function \(f\).

The finite-prime part of the argument involves fibers between which
no edge can pass. These fibers are \emph{edge-separated blocks};
they need not be connected components. The resulting formula is an
exact characterization, not an evaluation of \(\alpha\) as a
particular numerical constant. In particular, the present argument
does not determine that value and does not completely solve Problem 302.

\section{Parametrization and notation}

\begin{lemma}\label{lem:param}
For every solution of \eqref{eq:reciprocal} in integers
\(a<b<c\), there is a unique triple
\[
 (k,u,v)\in\N^3,\qquad u<v,\qquad (u,v)=1,
\]
for which \eqref{eq:param-intro} holds. Conversely, every such
triple gives a solution with \(a<b<c\).
\end{lemma}

\begin{proof}
Write \(d=(b,c)\), \(b=du\), and \(c=dv\), where \(u<v\) and
\((u,v)=1\). Equation \eqref{eq:reciprocal} gives
\[
 a=\frac{duv}{u+v}.
\]
Since \((uv,u+v)=1\), integrality implies \(u+v\mid d\).
Putting \(k=d/(u+v)\) proves the representation and its uniqueness.
Conversely, substitution verifies the equation, while
\[
 b-a=ku^2>0,\qquad c-b=k(v-u)(u+v)>0.
\]
\end{proof}

We regard a solution as the three-element edge \(\{a,b,c\}\).
The full hypergraph on \([N]\) has all such edges, and the truncated
hypergraph has those with \(v\le R\). A \emph{vertex cover} of a
specified family of edges is a subset of the vertex set meeting
every edge in that family.

For a real number \(x\ge2\) and \(n\in\N\), write
\[
 \Omega_x(n)=\sum_{\substack{p\le x\\p\ {\rm prime}}}v_p(n),
 \qquad
 \omega(n)=\#\{p:p\mid n\}.
\]
Here \(v_p(n)\) is the exponent of \(p\) in \(n\). Thus
\(\Omega_x\) counts prime factors with multiplicity. It is
completely additive:
\[
 \Omega_x(mn)=\Omega_x(m)+\Omega_x(n)
 \qquad(m,n\in\N).
\]
We use \(\mu\), \(\varphi\), and \(\pi\) for the M\"obius function,
Euler's totient function, and the prime-counting function,
respectively. A subscript indicates the allowed dependence of an
implied constant.

\section{A uniform Selberg marked-prime lemma}

We first record the elementary prime estimates needed in the sieve
argument.

\begin{lemma}\label{lem:prime-estimates}
Uniformly for every real \(x\ge3\),
\begin{equation}\label{eq:prime-estimates}
 \sum_{p\le x}\frac{\log p}{p}\ll\log x,
 \qquad
 \sum_{p\le x}\frac1p=\log\log x+O(1).
\end{equation}
The constants are effective and absolute.
\end{lemma}

\begin{proof}
The usual central-binomial-coefficient argument gives
\(\vartheta(x)=\sum_{p\le x}\log p\ll x\), with an effective
constant: for an integer \(n\ge1\),
\[
 \vartheta(2n)-\vartheta(n)
 \le \log\binom{2n}{n}\le 2n\log2,
\]
and summation over powers of two suffices. Partial summation gives
the first estimate in \eqref{eq:prime-estimates}.

For the upper bound in the second estimate, put
\(\sigma=1+1/\log x\). Then
\[
 \sum_{p\le x}\left(\frac1p-\frac1{p^\sigma}\right)
 \le \frac1{\log x}\sum_{p\le x}\frac{\log p}{p}\ll1.
\]
The Euler product and the integral bound for the zeta series give
\[
 \sum_{p\le x}p^{-\sigma}
 \le\log\zeta(\sigma)
 \le\log(1+\log x)=\log\log x+O(1).
\]
For the reverse inequality,
\[
 \prod_{p\le x}(1-p^{-1})^{-1}
 \ge \sum_{n\le x}\frac1n\ge\log x.
\]
Taking logarithms and using the absolute convergence of
\(\sum_p p^{-2}\) gives
\[
 \log\log x\le\sum_{p\le x}\frac1p+O(1).
\]
All bounds in this argument are effective.
\end{proof}

For primes \(p\ge7\), define
\[
 g_1(p)=\frac1p,\qquad
 g_2(p)=\frac{3p-2}{p^2}.
\]
Both satisfy
\begin{equation}\label{eq:g-bounds}
 0<g_i(p)<\frac12,\qquad g_i(p)\le\frac3p.
\end{equation}

\begin{lemma}[Uniform Selberg bound for arbitrary marked sets]
\label{lem:selberg-marked}
There exist effective absolute constants \(\theta>0\) and
\(C_S>0\) with the following properties.

Let \(x\ge3\) be real, let \(\mathcal P\) be any subset of the primes
in \([7,x]\), and put \(P=\prod_{p\in\mathcal P}p\), with the empty
product interpreted as \(1\).
\begin{enumerate}
\item For every real \(K\ge1\) satisfying \(x\le K^\theta\),
\begin{equation}\label{eq:marked-one}
 \#\{k\in\N:k\le K,\ (k,P)=1\}
 \le C_S K\prod_{p\in\mathcal P}(1-g_1(p)).
\end{equation}
\item For every real \(V\ge1\) satisfying \(x\le V^\theta\),
\begin{align}
 &\#\{(u,v)\in\N^2:1\le u\le2V,\ V\le v\le2V,\
                   (uv(u+v),P)=1\} \notag\\
 &\hspace{35mm}\le
 C_S V^2\prod_{p\in\mathcal P}(1-g_2(p)).
 \label{eq:marked-two}
\end{align}
\end{enumerate}
In particular, neither constant depends on the marked set
\(\mathcal P\).
\end{lemma}

\begin{proof}
Fix either local density \(g=g_i\), extend it multiplicatively to
squarefree divisors of \(P\), and set
\[
 h(p)=\frac{g(p)}{1-g(p)},\qquad
 h(d)=\prod_{p\mid d}h(p),\qquad
 G(D)=\sum_{\substack{d\mid P\\d\le D}}h(d)
 \quad(D\ge1).
\]

\paragraph{The Selberg quadratic form.}
There are real coefficients \(\lambda_d\), supported on
\(d\mid P\), \(d\le D\), such that \(\lambda_1=1\) and
\begin{equation}\label{eq:selberg-minimum}
 \sum_{d,e}\lambda_d\lambda_e g(\lcm(d,e))
 =\frac1{G(D)}.
\end{equation}
Indeed, put
\[
 y_r=\sum_{\substack{d\mid P\\r\mid d}}\lambda_d g(d).
\]
For squarefree \(d,e\),
\[
 g(\lcm(d,e))
 =g(d)g(e)\sum_{r\mid(d,e)}\frac1{h(r)}.
\]
Consequently,
\begin{equation}\label{eq:diagonalization}
 \sum_{d,e}\lambda_d\lambda_e g(\lcm(d,e))
 =\sum_{\substack{r\mid P\\r\le D}}\frac{y_r^2}{h(r)},
 \qquad
 \lambda_1=\sum_{\substack{r\mid P\\r\le D}}\mu(r)y_r.
\end{equation}
Cauchy--Schwarz gives the lower bound \(1/G(D)\) when
\(\lambda_1=1\), and equality is attained by
\[
 y_r=\frac{\mu(r)h(r)}{G(D)}
 \quad(r\mid P,\ r\le D),
 \qquad y_r=0\quad(r>D).
\]
M\"obius inversion gives the corresponding coefficients:
\begin{equation}\label{eq:lambda-formula}
 \lambda_d=
 \frac{\mu(d)}{G(D)}
 \prod_{p\mid d}(1-g(p))^{-1}
 \sum_{\substack{m\le D/d\\m\mid P\\(m,d)=1}}h(m)
 \quad(d\mid P,\ d\le D).
\end{equation}
This also proves directly that \(\lambda_1=1\). By positivity of
\(h\) and \eqref{eq:g-bounds},
\begin{equation}\label{eq:lambda-size}
 |\lambda_d|
 \le\prod_{p\mid d}(1-g(p))^{-1}
 \le2^{\omega(d)}\le d.
\end{equation}

\paragraph{A lower bound for \(G(D)\), uniform in \(\mathcal P\).}
Let
\[
 W=\sum_{d\mid P}h(d)
   =\prod_{p\in\mathcal P}(1-g(p))^{-1}.
\]
Give the divisors \(d\mid P\) the probability distribution
\(h(d)/W\). Under this distribution the prime factors of \(d\)
are independently selected, and \(p\) is selected with probability
\(g(p)\). Hence, by Lemma~\ref{lem:prime-estimates},
\[
 \mathbb E(\log d)
 =\sum_{p\in\mathcal P}g(p)\log p
 \le3\sum_{p\le x}\frac{\log p}{p}
 \le C_0\log x
\]
for an effective absolute \(C_0\).

Choose a fixed effective real \(A>\max(1,2C_0)\), and put
\(D=x^A\). Markov's inequality yields
\[
 \mathbb P(d>D)\le\frac{C_0}{A}<\frac12.
\]
Therefore
\begin{equation}\label{eq:G-uniform}
 G(D)\ge\frac12W,\qquad
 \frac1{G(D)}
 \le2\prod_{p\in\mathcal P}(1-g(p)).
\end{equation}
This is the required uniformity with respect to the marked set.
Fix from now on
\begin{equation}\label{eq:theta}
 \theta=\frac1{20A}.
\end{equation}

\paragraph{The one-dimensional remainder.}
For every positive integer \(m\), the Selberg majorant is
\begin{equation}\label{eq:selberg-majorant}
 \ind_{\{(m,P)=1\}}
 \le
 \left(\sum_{\substack{d\mid P\\d\le D\\d\mid m}}\lambda_d\right)^2.
\end{equation}
Apply this with \(m=k\), use \(g=g_1\), and use
\[
 \#\{k\in\N:k\le K,\ l\mid k\}=\frac Kl+O(1).
\]
Equations \eqref{eq:selberg-minimum} and
\eqref{eq:lambda-size} give
\begin{equation}\label{eq:one-remainder}
 \#\{k\le K:(k,P)=1\}
 \le\frac K{G(D)}+O(D^4),
\end{equation}
because \(\sum_{d\le D}|\lambda_d|\ll D^2\).

If \(x\le K^\theta\), then \(D\le K^{1/20}\), so the remainder
is \(O(K^{1/5})\). Moreover,
\[
 \prod_{p\in\mathcal P}(1-g_1(p))
 \ge\prod_{7\le p\le x}(1-p^{-1})
 \gg(\log x)^{-1}.
\]
The last estimate follows from Lemma~\ref{lem:prime-estimates}.
Since \(\log x\le\theta\log K\), the error in
\eqref{eq:one-remainder} is absorbed by a constant multiple of
the right side of \eqref{eq:marked-one}: for all sufficiently large
\(K\), \(K^{1/5}\log x\le K\), and the remaining bounded range is
absorbed by enlarging the absolute constant. This proves
\eqref{eq:marked-one} with an absolute constant.

\paragraph{The two-dimensional remainder.}
Over \(\mathbb F_p^2\), the union of the three lines
\[
 u=0,\qquad v=0,\qquad u+v=0
\]
has \(3p-2\) points. Thus, for squarefree \(l\mid P\), the number
of residue pairs satisfying \(l\mid uv(u+v)\) is
\[
 \varrho(l)=\prod_{p\mid l}(3p-2),
 \qquad \frac{\varrho(l)}{l^2}=g_2(l).
\]
Counting each residue pair in the rectangle gives
\begin{align}
 &\#\{(u,v)\in\N^2:1\le u\le2V,\ V\le v\le2V,\
                     l\mid uv(u+v)\} \notag\\
 &\hspace{15mm}
 =2V^2\frac{\varrho(l)}{l^2}
  +O\left(\varrho(l)\left(\frac Vl+1\right)\right).
 \label{eq:lattice-remainder}
\end{align}
The displayed error also includes the effects of nonintegral
endpoints and closed boundaries, so the nominal area \(2V^2\) is
valid for real \(V\).

Apply \eqref{eq:selberg-majorant} to \(m=uv(u+v)\) with
\(g=g_2\). Since \(\varrho(l)\le l^2\) and
\(\lcm(d,e)\le D^2\), the total error after expansion is
\[
 \ll
 \sum_{d,e\le D}|\lambda_d\lambda_e|
 \left(V\lcm(d,e)+\lcm(d,e)^2\right)
 \ll VD^6+D^8.
\]
It follows that the left side of \eqref{eq:marked-two} is at most
\begin{equation}\label{eq:two-remainder}
 \frac{2V^2}{G(D)}+O(VD^6+D^8).
\end{equation}
If \(x\le V^\theta\), then \(D\le V^{1/20}\), and the error is
\[
 O(V^{13/10}+V^{2/5}).
\]
Finally,
\[
 \prod_{p\in\mathcal P}(1-g_2(p))
 \ge\prod_{7\le p\le x}(1-g_2(p))
 \gg(\log x)^{-3},
\]
because \(g_2(p)=3/p+O(p^{-2})\). The error in
\eqref{eq:two-remainder} is therefore absorbed into
\(O(V^2\prod_{p\in\mathcal P}(1-g_2(p)))\): for all sufficiently
large \(V\), \((V^{13/10}+V^{2/5})(\log x)^3\le V^2\), and the
remaining bounded range is absorbed by the absolute constant.
Together with \eqref{eq:G-uniform}, this proves
\eqref{eq:marked-two}. All constants used above are independent
of \(\mathcal P\).
\end{proof}

\begin{corollary}[Weighted sieve estimates]\label{cor:weighted}
Let \(\theta>0\) be as in Lemma~\ref{lem:selberg-marked}.
For every fixed real \(q\in(0,1)\), there is an effective constant
\(C_q>0\) such that
\begin{equation}\label{eq:weighted-one}
 \sum_{\substack{k\in\N\\k\le K}}q^{\Omega_x(k)}
 \le C_qK(\log x)^{q-1}
 \qquad(K\ge1,\ 3\le x\le K^\theta),
\end{equation}
and
\begin{equation}\label{eq:weighted-two}
 \sum_{\substack{(u,v)\in\N^2\\1\le u\le2V\\V\le v\le2V}}
 q^{\Omega_x(uv(u+v))}
 \le C_qV^2(\log x)^{3(q-1)}
 \qquad(V\ge1,\ 3\le x\le V^\theta).
\end{equation}
Here \(K,V,x\) may be real. No coprimality condition on \(u,v\)
is imposed.
\end{corollary}

\begin{proof}
Independently mark each prime in \([7,x]\) with probability \(1-q\),
and let \(\mathcal P\) be the random marked set. For every
\(m\in\N\),
\[
 q^{\Omega_x(m)}
 \le q^{\#\{p\in[7,x]:p\mid m\}}
 =\mathbb P\bigl((m,\prod_{p\in\mathcal P}p)=1\bigr).
\]
The inequality uses \(0<q<1\); omitting the primes \(2,3,5\)
and ignoring multiplicities can only increase the weight.

Average \eqref{eq:marked-one} or \eqref{eq:marked-two}. Their
constants are independent of \(\mathcal P\), and independence
of the marks gives
\[
 \mathbb E\prod_{p\in\mathcal P}(1-g_i(p))
 =\prod_{7\le p\le x}\bigl(1-(1-q)g_i(p)\bigr).
\]
Writing \(\kappa_1=1\), \(\kappa_2=3\), we have
\(g_i(p)=\kappa_i/p+O(p^{-2})\). Lemma~\ref{lem:prime-estimates}
therefore gives
\[
 \prod_{7\le p\le x}\bigl(1-(1-q)g_i(p)\bigr)
 \asymp_q(\log x)^{-\kappa_i(1-q)}.
\]
This proves both assertions.
\end{proof}

\section{Divisor intervals and a simultaneous exceptional set}

\subsection{The input from Ford's theorem}

For real \(X\ge1\) and \(1<y<z\), let
\[
 H(X,y,z)=
 \#\{n\in\N:n\le X,\ \text{there exists }d\mid n
                         \text{ with }y<d\le z\}.
\]
Put
\begin{equation}\label{eq:delta}
 \delta=
 1-\frac{1+\log\log2}{\log2}
 =0.08607\ldots.
\end{equation}

\begin{theorem}[Ford, upper-bound specialization]\label{thm:ford}
There exist absolute constants \(C_F>0\) and \(y_F\ge3\)
such that, for all real \(X,y,z\) satisfying
\[
 y\ge y_F,\qquad 2y\le z\le\min(y^2,\sqrt X),
\]
one has
\begin{equation}\label{eq:ford}
 H(X,y,z)
 \le C_F X
       u^\delta\bigl(\log(2/u)\bigr)^{-3/2},
 \qquad
 u=\frac{\log(z/y)}{\log y}.
\end{equation}
\end{theorem}

This is the indicated subrange of Kevin Ford,
\emph{The distribution of integers with a divisor in a given
interval}, \emph{Annals of Mathematics} \textbf{168} (2008),
367--433, Theorem~1(v), p.~371, DOI
\href{https://doi.org/10.4007/annals.2008.168.367}
{10.4007/annals.2008.168.367}.
Only the upper bound \eqref{eq:ford}, with its uniform constants, is
used here. In particular, the application is not
restricted to intervals with fixed ratio \(z/y\).
In Ford's notation, the middle case of Theorem~1(v) assumes
\(X>100000\), \(100\le y\le z-1\), \(y\le\sqrt X\), and
\(2y\le z\le y^2\). Enlarging \(y_F\) absorbs the fixed lower
thresholds on \(X\) and \(y\), while our hypotheses retain
\(y\le z/2\le z-1\), \(z\le\sqrt X\), and the same middle-case
range. Thus the displayed specialization is a direct restriction of
Ford's theorem.

\begin{corollary}\label{cor:ford-window}
There are absolute constants \(M_F\ge1\) and \(C_D>0\)
such that for every \(N\in\N\) and every real
\(0<\eta\le1/16\) satisfying \(\eta\log N\ge M_F\), the set
\[
 \mathcal D_\eta(N)=
 \left\{n\in[N]:
   \text{there exists }d\mid n,\quad
   N^{1/2-2\eta}\le d\le N^{1/2}
 \right\}
\]
satisfies
\begin{equation}\label{eq:ford-window}
 |\mathcal D_\eta(N)|\le C_DN\eta^\delta.
\end{equation}
\end{corollary}

\begin{proof}
Set \(y=N^{1/2-2\eta}\), \(z=N^{1/2}\). Then
\[
 \frac zy=N^{2\eta},\qquad
 z\le y^2,\qquad
 u=\frac{4\eta}{1-4\eta}.
\]
Here \(z\le y^2\) follows from \(\eta\le1/16\). Choose \(M_F\)
large enough that \(\eta\log N\ge M_F\) implies both \(2y\le z\)
and \(y\ge y_F\). Such a choice is absolute: indeed,
\[
 \log y\ge\frac38\log N\ge6M_F.
\]
Moreover,
\[
 4\eta\le u\le\frac{16}3\eta\le\frac13.
\]
Theorem~\ref{thm:ford} now gives
\(H(N,y,z)\ll N\eta^\delta\), after discarding the favorable
logarithmic factor.

The definition of \(\mathcal D_\eta(N)\) also permits \(d=y\).
If \(y\) is an integer, these additional integers number at most
\(N/y\le N^{5/8}\). Since \(M_F\ge1\), the hypothesis gives
\(\eta\ge1/\log N\), and
\[
 N^{5/8}\ll N(\log N)^{-\delta}\le N\eta^\delta.
\]
If \(y\) is not an integer, there is no endpoint contribution.
This proves \eqref{eq:ford-window}.
\end{proof}

\subsection{Simultaneous control of small prime factors}

Fix the absolute parameters
\begin{equation}\label{eq:fixed-parameters}
 \eps=\frac2{25}=0.08,\qquad
 t=\frac{26}{25}=1.04,\qquad
 c_*=(1+\eps)\log t-(t-1)>0.
\end{equation}

\begin{lemma}\label{lem:exceptional}
For real \(T\ge3\) and \(N\in\N\), define
\[
 \mathcal B_T(N)=
 \left\{n\in[N]:
   \Omega_x(n)>(1+\eps)\log\log x
   \text{ for some real }x\ge T
 \right\}.
\]
There is an effective absolute constant \(C_B>0\) such that
\begin{equation}\label{eq:exceptional}
 |\mathcal B_T(N)|
 \le C_BN(\log T)^{-c_*}
 \qquad(T\ge3,\ N\in\N).
\end{equation}
\end{lemma}

\begin{proof}
We first prove, uniformly for real \(x\ge3\) and \(N\in\N\),
\begin{equation}\label{eq:positive-moment}
 \frac1N\sum_{n\le N}t^{\Omega_x(n)}
 \ll_t(\log x)^{t-1}.
\end{equation}
Define a nonnegative multiplicative function \(\gamma_x\) by
\[
 \gamma_x(p^a)=
 \begin{cases}
  (t-1)t^{a-1},&p\le x,\ a\ge1,\\
  0,&p>x,\ a\ge1.
 \end{cases}
\]
Then
\[
 t^{\Omega_x(n)}=\sum_{d\mid n}\gamma_x(d).
\]
Because \(1<t<2\), every local series below converges, and
positivity gives
\begin{align*}
 \frac1N\sum_{n\le N}t^{\Omega_x(n)}
 &\le\sum_{d\ge1}\frac{\gamma_x(d)}d\\
 &=\prod_{p\le x}\frac{1-1/p}{1-t/p}
 \ll_t(\log x)^{t-1}.
\end{align*}
For the last estimate, take logarithms and apply
Lemma~\ref{lem:prime-estimates}. This proves
\eqref{eq:positive-moment}, even when \(x>N\).

Let \(x_j=\exp(e^j)\) for integers \(j\ge0\), and put
\(j_0=\lfloor\log\log T\rfloor\), which is nonnegative.
If \(x\in[x_j,x_{j+1}]\) witnesses membership in
\(\mathcal B_T(N)\), then
\[
 \Omega_{x_{j+1}}(n)>(1+\eps)j.
\]
By Markov's inequality and \eqref{eq:positive-moment}, the number
of integers with this property is at most
\[
 \ll N t^{-(1+\eps)j}
          (\log x_{j+1})^{t-1}
 =O\left(N e^{-c_*j}\right).
\]
Summing over all \(j\ge j_0\) gives
\[
 |\mathcal B_T(N)|
 \ll N e^{-c_*j_0}
 \ll N(\log T)^{-c_*}.
\]
The infinite grid is legitimate because the moment estimate is
uniform in both \(N\) and \(x\).
\end{proof}

\section{A uniform cover for the primitive tail}

Fix
\begin{equation}\label{eq:z-beta}
 z=\frac{\sqrt5-1}{2},\qquad
 \beta=4-3z^2-z^3+3(1+\eps)\log z.
\end{equation}
The following rationally certified numerical inequalities will suffice:
\begin{equation}\label{eq:numeric}
 \delta>0.086,\qquad c_*>0.0023,\qquad \beta>1.058.
\end{equation}
In particular, \(\beta>1\). For completeness, the following explicit
intervals certify the displayed inequalities:
\[
\begin{gathered}
 2.236067977<\sqrt5<2.236067978,\\
 0.6931471805<\log2<0.6931471806,\\
 -0.3665129210<\log\log2<-0.3665129203,\\
 -0.481211826<\log z<-0.481211824,\\
 0.0392207131<\log(26/25)<0.0392207132.
\end{gathered}
\]
Substitution with outward rounding gives
\[
 0.08607<\delta<0.08608,\qquad
 0.00235<c_*<0.00237,\qquad
 1.0588<\beta<1.0590.
\]
Each logarithmic interval follows by taking \(M=20\) in
\[
 \log a=
 2\sum_{m=0}^{M-1}\frac{w^{2m+1}}{2m+1}
 +\mathcal R_M,\qquad
 w=\frac{a-1}{a+1},\qquad
 |\mathcal R_M|
 \le\frac{2|w|^{2M+1}}{(2M+1)(1-w^2)}
 \quad(a>0).
\]
For \(\log\log2\), first enclose \(\log2\) by the same formula
and then apply it once more with interval arithmetic. The square-root
interval follows by squaring its rational endpoints.

\begin{proposition}[Parameterized tail cover]\label{prop:param-cover}
Let \(\theta\) be as in Corollary~\ref{cor:weighted} and
\(M_F\) as in Corollary~\ref{cor:ford-window}. There is an absolute
constant \(C_1>0\) with the following property.

For every integer \(R\ge2\), every \(N\in\N\), and every real
\(0<\eta\le1/16\) satisfying
\begin{equation}\label{eq:cover-hypotheses}
 T:=R^{2\eta\theta}\ge3,\qquad
 \eta\log N\ge M_F,\qquad
 N\ge\eta^{-1/\eta},
\end{equation}
there exists a set \(\mathcal T_{R,\eta}(N)\subseteq[N]\)
meeting every solution \eqref{eq:param-intro} with \(v>R\) and
\(c\le N\), such that
\begin{align}
 \frac{|\mathcal T_{R,\eta}(N)|}{N}
 \le C_1\bigl(
   \eta+\eta^\delta
   +(\eta\log R)^{-c_*}
   +\eta^{-\beta}(\log R)^{1-\beta}
 \bigr).
 \label{eq:param-cover}
\end{align}
The constant \(C_1\) is independent of \(R,N,\eta\).
\end{proposition}

\begin{proof}
Begin by deleting
\[
 \mathcal T_0=
 \{n\in[N]:n\le\eta N\}
 \ \cup\ \mathcal D_\eta(N)
 \ \cup\ \mathcal B_T(N).
\]
Corollary~\ref{cor:ford-window} and
Lemma~\ref{lem:exceptional} give
\begin{equation}\label{eq:initial-cost}
 \frac{|\mathcal T_0|}{N}
 \ll \eta+\eta^\delta+(\eta\log R)^{-c_*},
\end{equation}
where \(\log T=2\eta\theta\log R\), and \(\theta\) is fixed.

\paragraph{Solutions with \(k<N^\eta\).}
Suppose that a solution survives deletion of the first two sets
in \(\mathcal T_0\), and that \(k<N^\eta\).
Since \(a>\eta N\) and
\[
 v^2<kv(u+v)=c\le N,
\]
we have
\[
 u=\frac a{kv}>\eta N^{1/2-\eta}.
\]
The hypothesis \(N\ge\eta^{-1/\eta}\) is equivalent to
\(\eta\ge N^{-\eta}\). Hence
\[
 N^{1/2-2\eta}<u<v<N^{1/2}.
\]
But \(u\mid a\), so \(a\in\mathcal D_\eta(N)\), a contradiction.
Thus no solution with \(k<N^\eta\) survives \(\mathcal T_0\).

\paragraph{Solutions with \(k\ge N^\eta\).}
Consider a dyadic range
\[
 V\le v<2V,\qquad V=2^jR,\quad j\in\mathbb Z_{\ge0},
\]
which contains at least one surviving solution with
\(k\ge N^\eta\). Put
\begin{equation}\label{eq:scale}
 K=\frac N{V^2},\qquad x=\min(V,K)^\theta.
\end{equation}
For any solution in this range,
\[
 c=kv(u+v)>kV^2,
\]
so \(k<K\). Nonemptiness therefore gives
\[
 K\ge N^\eta,\qquad N>V^2,\qquad
 K\ge V^{2\eta}.
\]
Since \(2\eta\le1\),
\begin{equation}\label{eq:x-range}
 x\ge V^{2\eta\theta}\ge R^{2\eta\theta}=T,
 \qquad x\le V^\theta,\qquad x\le K^\theta.
\end{equation}
Both weighted sieve estimates apply at this same value of \(x\).

All three vertices of a surviving solution lie outside
\(\mathcal B_T(N)\). By complete additivity of \(\Omega_x\),
\begin{align}
 S&:=\Omega_x(a)+\Omega_x(b)+\Omega_x(c)\notag\\
  &=3\Omega_x(k)+2\Omega_x(uv(u+v))\notag\\
  &\le 3(1+\eps)\log\log x.
 \label{eq:normal-edge}
\end{align}
This identity does not require any coprimality between the
displayed factors.

If \(0<z<1\) and \(S\le A\), then \(1\le z^{S-A}\).
Consequently, the number of surviving solutions in the range
under consideration is at most
\begin{align}
 &(\log x)^{3(1+\eps)\log(1/z)}
 \left(\sum_{k\le K}(z^3)^{\Omega_x(k)}\right)
 \left(
 \sum_{\substack{1\le u\le2V\\V\le v\le2V}}
 (z^2)^{\Omega_x(uv(u+v))}
 \right).
 \label{eq:weighted-edge-count}
\end{align}
We have enlarged the parameter range, dropping both \(u<v\)
and \((u,v)=1\), which is valid for this upper bound.

Apply Corollary~\ref{cor:weighted} with \(q=z^3\) and \(q=z^2\).
The expression in \eqref{eq:weighted-edge-count} is
\begin{align}
 &\ll KV^2
 (\log x)^{3(1+\eps)\log(1/z)+(z^3-1)+3(z^2-1)}
 \notag\\
 &=N(\log x)^{-\beta}
 \ll N\eta^{-\beta}(\log V)^{-\beta}.
 \label{eq:dyadic-bound}
\end{align}
The last inequality follows from \eqref{eq:x-range}. Its implied
constant is absolute because \(\theta,z,\eps\) are fixed.

The relevant dyadic ranges are finite, but we may sum over all
\(j\ge0\). Since \(\beta>1\),
\[
 \sum_{j\ge0}(\log R+j\log2)^{-\beta}
 \ll(\log R)^{1-\beta}
 \qquad(R\ge2).
\]
Thus the number of all surviving tail edges is at most
\begin{equation}\label{eq:surviving-edges}
 \ll N\eta^{-\beta}(\log R)^{1-\beta}.
\end{equation}
Choose one vertex from every such edge and adjoin all chosen
vertices to \(\mathcal T_0\). The union meets every tail edge,
and its additional size is at most the number in
\eqref{eq:surviving-edges}. Combining this with
\eqref{eq:initial-cost} proves \eqref{eq:param-cover}.
\end{proof}

\begin{theorem}[Uniform tail cover]\label{thm:tail}
There exist absolute constants \(C>0\) and
\(R_0\in\N\), \(R_0\ge2\), such that for every integer \(R\ge R_0\)
and every \(N\in\N\), there is a set
\(\mathcal T_R(N)\subseteq[N]\) meeting every solution
\eqref{eq:param-intro} with \(v>R\) and \(c\le N\), and satisfying
\begin{equation}\label{eq:tail-cover}
 |\mathcal T_R(N)|\le CN(\log R)^{-1/1000}.
\end{equation}
\end{theorem}

\begin{proof}
If \(N\le R^2\), no tail edge exists, because every such edge
satisfies \(c>v^2>R^2\). In this case take
\(\mathcal T_R(N)=\varnothing\).

Suppose \(N>R^2\), put \(L=\log R\), and choose
\[
 \eta=L^{-1/40}.
\]
For all sufficiently large \(R\), uniformly over
\(N>R^2\), the hypotheses of Proposition~\ref{prop:param-cover}
hold. Indeed,
\[
 \eta\le\frac1{16},\qquad
 \log T=2\theta L^{39/40}\longrightarrow\infty,
 \qquad
 \eta\log N>2L^{39/40}\longrightarrow\infty,
\]
and
\[
 \log(\eta^{-1/\eta})
 =\frac1{40}L^{1/40}\log L=o(L)<\log N.
\]
The thresholds can be chosen to hold for every larger \(R\).

The four terms in \eqref{eq:param-cover} now become
\[
 L^{-1/40},\qquad
 L^{-\delta/40},\qquad
 L^{-(39/40)c_*},\qquad
 L^{-(\beta-1-\beta/40)}.
\]
By \eqref{eq:numeric}, every displayed decay exponent is greater
than \(1/1000\). For the last one, for example,
\[
 \beta-1-\frac{\beta}{40}
 =\frac{39}{40}\beta-1
 >\frac{39}{40}(1.058)-1
 >\frac1{1000}.
\]
Proposition~\ref{prop:param-cover} proves
\eqref{eq:tail-cover} after increasing an absolute constant.
\end{proof}

\begin{corollary}\label{cor:extremal-comparison}
With \(C,R_0\) as in Theorem~\ref{thm:tail}, for every integer
\(R\ge R_0\) and every \(N\in\N\),
\[
 0\le f_R(N)-f(N)\le CN(\log R)^{-1/1000}.
\]
\end{corollary}

\begin{proof}
Let \(A\subseteq[N]\) be a maximum independent set in the
truncated hypergraph. Then \(A\setminus\mathcal T_R(N)\) contains
neither a truncated edge nor a tail edge. Therefore
\[
 f(N)\ge |A\setminus\mathcal T_R(N)|
 \ge f_R(N)-|\mathcal T_R(N)|.
\]
The reverse comparison \(f(N)\le f_R(N)\) follows directly from
the definitions.
\end{proof}

\begin{remark}
The set \(\mathcal T_R(N)\) is a cover for the entire tail
hypergraph, not merely for the tail edges in a selected
independent set. The proof gives an upper bound for
\(f_R(N)-f(N)\); it does not assert equality with the minimum
size of a tail cover.
\end{remark}

\section{Finite-prime decomposition}

Fix an integer \(R\ge2\), and define
\begin{equation}\label{eq:Q}
 Q_R=\prod_{p\le2R}p,\qquad
 \rho_R=\frac{\varphi(Q_R)}{Q_R},\qquad
 D_R=2^{\pi(2R)}.
\end{equation}
Call a positive integer \(Q_R\)-smooth if all its prime factors
divide \(Q_R\), with \(1\) included. List these integers as
\[
 1=s_{R,1}<s_{R,2}<\cdots,
\]
and, for every real \(y\ge0\), put
\[
 \Psi_R(y)=\#\{j:s_{R,j}\le y\},\qquad
 C_R(y)=\#\{r\in\N:r\le y,\ (r,Q_R)=1\}.
\]

\begin{definition}\label{def:finite-variational}
For each integer \(j\ge1\), let \(\mathcal G_{R,j}\) be the
hypergraph on
\(\{s_{R,1},\ldots,s_{R,j}\}\) whose edges are the solutions
\eqref{eq:param-intro} with \(v\le R\) and with all three vertices
in the displayed set. Let \(h_{R,j}\) be its
independence number, and put \(h_{R,0}=0\). Equivalently,
\begin{equation}\label{eq:finite-max}
 h_{R,j}=
 \max\left\{
 \sum_{i=1}^j\xi_i:
 \begin{array}{l}
 \xi_i\in\{0,1\}\quad(1\le i\le j),\\
 \xi_{i_1}+\xi_{i_2}+\xi_{i_3}\le2
 \text{ for every edge }
 \{s_{R,i_1},s_{R,i_2},s_{R,i_3}\}
 \end{array}
 \right\}.
\end{equation}
Define
\[
 d_{R,j}=h_{R,j}-h_{R,j-1}\qquad(j\ge1).
\]
\end{definition}

Because \(\mathcal G_{R,j-1}\) is an induced subhypergraph of
\(\mathcal G_{R,j}\), an independent set in the former remains
independent in the latter. Conversely, removing the new vertex
from any independent set loses at most one vertex. Hence
\begin{equation}\label{eq:increments}
 d_{R,j}\in\{0,1\},\qquad
 h_{R,m}=\sum_{j=1}^m d_{R,j}\quad(m\ge0).
\end{equation}
No choice of nested maximizing independent sets is involved.

\begin{proposition}[Edge-separated blocks and exact decomposition]
\label{prop:blocks}
For every integer \(R\ge2\) and every \(N\in\N\),
\begin{align}
 f_R(N)
 &=\sum_{\substack{r\le N\\(r,Q_R)=1}}
       h_{R,\Psi_R(N/r)}
 \label{eq:block-sum}\\
 &=\sum_{j\ge1}d_{R,j}\,C_R(N/s_{R,j}).
 \label{eq:exact-count}
\end{align}
The series in \eqref{eq:exact-count} has only finitely many
nonzero terms.
\end{proposition}

\begin{proof}
Every integer \(n\in\N\) has a unique factorization
\[
 n=rs,\qquad (r,Q_R)=1,\qquad s\text{ is }Q_R\text{-smooth}.
\]
For each \(r\le N\) with \((r,Q_R)=1\), define the fiber
\[
 V_r(N)=\{rs\le N:s\text{ is }Q_R\text{-smooth}\}.
\]
These fibers partition \([N]\).

In a truncated edge, \(u<v\le R\), and every prime dividing
\(uv(u+v)\) is at most \(2R\). Write the scaling parameter as
\(k=rt\), where \((r,Q_R)=1\) and \(t\) is \(Q_R\)-smooth.
All three vertices then have the same factor \(r\) supported on
primes outside \(Q_R\). Thus no truncated edge crosses between
different fibers.

Division by \(r\) identifies the induced hypergraph on \(V_r(N)\)
with \(\mathcal G_{R,\Psi_R(N/r)}\). This identification holds in
both directions: multiplication by \(r\) changes \(k\) to \(rk\)
without changing \(u,v\), while division of an edge in the fiber
changes \(k=rt\) to \(t\). Independence numbers therefore add
over these edge-separated blocks, proving \eqref{eq:block-sum}.

Now use \eqref{eq:increments} and rearrange a finite sum:
\begin{align*}
 f_R(N)
 &=\sum_{\substack{r\le N\\(r,Q_R)=1}}
       \sum_{s_{R,j}\le N/r}d_{R,j}\\
 &=\sum_{s_{R,j}\le N}
       d_{R,j}\,\#\{r\le N/s_{R,j}:(r,Q_R)=1\}.
\end{align*}
Terms with \(s_{R,j}>N\) vanish, giving \eqref{eq:exact-count}.
\end{proof}

\begin{remark}\label{rem:not-components}
The fibers \(V_r(N)\) are edge-separated blocks, not necessarily
connected components. For example, when \(R=2\) and \(N=6\),
one has \(Q_R=6\) and
\[
 V_1(6)=\{1,2,3,4,6\}.
\]
Its only edge is \(\{2,3,6\}\); the vertices \(1\) and \(4\) are
isolated. Independence-number additivity requires only that no
edge cross between fibers.
\end{remark}

\begin{proposition}[Truncated limits and their error bounds]
\label{prop:truncated-limit}
For every integer \(R\ge2\), the limit \(\alpha_R\) exists and is
given by the absolutely convergent series
\begin{equation}\label{eq:alphaR-series}
 \alpha_R=\rho_R\sum_{j\ge1}\frac{d_{R,j}}{s_{R,j}}.
\end{equation}
Define
\begin{equation}\label{eq:B}
 B_R=\prod_{p\le2R}(1-p^{-1/2})^{-1}.
\end{equation}
Then for every \(N\in\N\) and every real \(X\ge1\),
\begin{equation}\label{eq:truncated-error}
 \left|\frac{f_R(N)}N-\alpha_R\right|
 \le\frac{D_R\Psi_R(X)}N+2B_RX^{-1/2}.
\end{equation}
In addition, if
\begin{equation}\label{eq:L}
 L_{R,X}=\rho_R
          \sum_{s_{R,j}\le X}\frac{d_{R,j}}{s_{R,j}},
\end{equation}
then
\begin{equation}\label{eq:finite-alphaR}
 L_{R,X}\le\alpha_R
 \le L_{R,X}+\rho_RB_RX^{-1/2}.
\end{equation}
\end{proposition}

\begin{proof}
Inclusion--exclusion gives, for every real \(y\ge0\),
\begin{equation}\label{eq:coprime-count}
 C_R(y)=\sum_{d\mid Q_R}\mu(d)\left\lfloor\frac yd\right\rfloor
       =\rho_Ry+E_R(y),\qquad |E_R(y)|\le D_R.
\end{equation}
Also,
\begin{equation}\label{eq:smooth-Euler}
 \sum_{j\ge1}\frac1{s_{R,j}}
 =\prod_{p\le2R}(1-p^{-1})^{-1}
 =\rho_R^{-1}<\infty.
\end{equation}
For every \(j\) and \(N\),
\[
 0\le\frac{d_{R,j}C_R(N/s_{R,j})}{N}
 \le\frac1{s_{R,j}}.
\]
For fixed \(j\), \eqref{eq:coprime-count} implies convergence of
this summand to \(\rho_Rd_{R,j}/s_{R,j}\) as \(N\to\infty\).
Dominated convergence applied to \eqref{eq:exact-count}, using
\eqref{eq:smooth-Euler}, proves \eqref{eq:alphaR-series}.
Here the finite sum in \eqref{eq:exact-count} is viewed as an infinite
series by setting \(C_R(N/s)=0\) whenever \(s>N\).
In particular, we do not sum the error in
\eqref{eq:coprime-count} over infinitely many indices.

For \(X\ge1\), the elementary bound
\begin{equation}\label{eq:rankin}
 \sum_{s_{R,j}>X}\frac1{s_{R,j}}
 \le X^{-1/2}\sum_{j\ge1}s_{R,j}^{-1/2}
 =B_RX^{-1/2}
\end{equation}
follows by summing the finite-prime Euler product at exponent
\(1/2\). Split both \eqref{eq:exact-count} and
\eqref{eq:alphaR-series} at \(X\). By
\eqref{eq:coprime-count}, the difference of their heads, after
normalization by \(N\), is at most \(D_R\Psi_R(X)/N\).
Their two tails together are at most
\[
 (1+\rho_R)\sum_{s_{R,j}>X}\frac1{s_{R,j}}
 \le2B_RX^{-1/2}.
\]
This proves \eqref{eq:truncated-error}, with no restriction
\(X\le N\). Finally, the summands of
\eqref{eq:alphaR-series} are nonnegative and \(d_{R,j}\le1\),
so \eqref{eq:rankin} also proves \eqref{eq:finite-alphaR}.
\end{proof}

\section{The limiting density and the variational formula}

\begin{proof}[Proof of Theorem~\ref{thm:main}, parts (1)--(3)]
Part (1) is Proposition~\ref{prop:truncated-limit}, and part (2)
is Corollary~\ref{cor:extremal-comparison}.
Since \(f_{R+1}(N)\le f_R(N)\), the numbers \(\alpha_R\) form a
nonincreasing sequence bounded below by \(0\). Put
\[
 a_*=\inf_{\substack{R\ge2\\R\in\N}}\alpha_R.
\]
For every fixed integer \(R\ge R_0\),
Corollary~\ref{cor:extremal-comparison} and the truncated limit
give
\[
 \alpha_R-C(\log R)^{-1/1000}
 \le\liminf_{N\to\infty}\frac{f(N)}N
 \le\limsup_{N\to\infty}\frac{f(N)}N
 \le\alpha_R.
\]
Letting \(R\to\infty\) shows that both middle quantities equal
\(a_*\). This proves existence of \(\alpha\) and the infimum
formula. Taking \(N\to\infty\) in the uniform comparison then
gives \eqref{eq:alpha-inf-main}.
\end{proof}

\begin{theorem}[Variational characterization and truncation bounds]
\label{thm:effective}
Use the definitions
\eqref{eq:Q}, \eqref{eq:finite-max}, \eqref{eq:B}, and
\eqref{eq:L}, and let \(C,R_0\) be as in
Theorem~\ref{thm:tail}. Then
\begin{equation}\label{eq:variational}
 \boxed{\displaystyle
 \alpha=
 \inf_{\substack{R\ge2\\R\in\N}}
 \left[
 \frac{\varphi(Q_R)}{Q_R}
 \sum_{j\ge1}
 \frac{h_{R,j}-h_{R,j-1}}{s_{R,j}}
 \right].}
\end{equation}
For every integer \(R\ge R_0\) and every real \(X\ge1\),
\begin{equation}\label{eq:certified-interval}
 \boxed{\displaystyle
 L_{R,X}-C(\log R)^{-1/1000}
 \le\alpha
 \le L_{R,X}+\rho_RB_RX^{-1/2}.}
\end{equation}
Furthermore, for every such \(R,X\) and every \(N\in\N\),
\begin{equation}\label{eq:full-effective}
 \boxed{\displaystyle
 \left|\frac{f(N)}N-\alpha\right|
 \le C(\log R)^{-1/1000}
       +\frac{D_R\Psi_R(X)}N
       +2B_RX^{-1/2}.}
\end{equation}
All finite optimization problems occurring in \(L_{R,X}\) are finite
and can be evaluated exactly. The constants \(C\) and \(R_0\) in
these bounds are not numerically extracted here.
\end{theorem}

\begin{proof}
Equation \eqref{eq:variational} is the combination of
\eqref{eq:alpha-inf-main} and \eqref{eq:alphaR-series}.
Put
\[
 \epsilon_R=C(\log R)^{-1/1000}\qquad(R\ge R_0).
\]
We have \(0\le\alpha_R-\alpha\le\epsilon_R\), while
\eqref{eq:finite-alphaR} gives
\[
 L_{R,X}\le\alpha_R
 \le L_{R,X}+\rho_RB_RX^{-1/2}.
\]
These inequalities imply \eqref{eq:certified-interval}.

For \eqref{eq:full-effective}, define
\[
 e_N=\frac{f_R(N)-f(N)}N,\qquad
 e_\infty=\alpha_R-\alpha.
\]
Both lie in \([0,\epsilon_R]\), and hence
\(|e_N-e_\infty|\le\epsilon_R\). Since
\[
 \frac{f(N)}N-\alpha
 =
 \left(\frac{f_R(N)}N-\alpha_R\right)
 -(e_N-e_\infty),
\]
Proposition~\ref{prop:truncated-limit} proves
\eqref{eq:full-effective}. The one-sided comparisons are why
only one copy of \(\epsilon_R\) is required.
\end{proof}

Theorem~\ref{thm:effective} proves part (4) of
Theorem~\ref{thm:main}.

\section{Scope of the conclusion}

The proof separates two issues. For a fixed cutoff \(R\),
finite-prime factorization gives an exact decomposition and a
convergent series for the truncated extremum. This fact alone
does not justify passage to the full problem. The uniform tail
cover supplies that missing comparison.

Its analytic mechanism is the strict inequality \(\beta>1\).
After integers with excessive small-prime-factor counts are
removed, the weighted count of tail edges is summable over
dyadic scales. The range where the scaling parameter is too
small for this simultaneous sieve application is covered
instead by Ford's divisor-interval estimate. All deletions
are uniform in the ambient endpoint \(N\).

The variational characterization is not computationally efficient:
its truncations may require extremely large cutoffs and finite
independence computations. No identification of \(\alpha\) as a
particular rational or other familiar constant is obtained here, and
no numerical improvement is claimed. Thus the existence of the
limiting density is established, but the present argument does not
completely solve Problem 302.

\begin{thebibliography}{9}
\sloppy

\bibitem{EG1980}
Paul Erd\H{o}s and Ronald L. Graham,
\emph{Old and New Problems and Results in Combinatorial Number Theory},
Monographies de L'Enseignement Math\'ematique, vol.~28,
Universit\'e de Gen\`eve, Geneva, 1980, Chapter~III.

\bibitem{EP302}
Thomas F. Bloom,
Erd\H{o}s Problem 302,
\emph{Erd\H{o}s Problems}, accessed 17 September 2026,
\url{https://www.erdosproblems.com/302}.

\bibitem{Khanukov2026}
Dmitry Khanukov,
Two-sided partial progress on Erd\H{o}s Problem 302,
corrected preprint, version 0.2.0, 2026.
Zenodo, DOI:
\href{https://doi.org/10.5281/zenodo.22737316}
{10.5281/zenodo.22737316}.

\bibitem{Ford2008}
Kevin Ford,
The distribution of integers with a divisor in a given interval,
\emph{Annals of Mathematics} (2) \textbf{168} (2008), no.~2,
367--433.
Theorem~1.
DOI:
\href{https://doi.org/10.4007/annals.2008.168.367}
{10.4007/annals.2008.168.367}.

\bibitem{HR}
Heini Halberstam and Hans-Egon Richert,
\emph{Sieve Methods},
London Mathematical Society Monographs, vol.~4,
Academic Press, London, 1974.

\end{thebibliography}

\end{document}